% =====================================================================
%  fisica-estatico.tex
%  Fallback para XeLaTeX (UaiTeX, Overleaf com XeLaTeX etc.).
%  GERADO por tools/gerar-fallback-xetex.lua -- nao edite a mao.
%
%  Para atualizar depois de mexer em lua/ ou dados/:
%      lua tools/gerar-fallback-xetex.lua
% =====================================================================
\makeatletter

\expandafter\newcommand\csname fisica@cte@G\endcsname{$G = \qty{6.67430(15)e-11}{\cubic\metre\per\kilogram\per\square\second}$}
\expandafter\newcommand\csname fisica@valor@G\endcsname{\qty{6.67430(15)e-11}{\cubic\metre\per\kilogram\per\square\second}}
\expandafter\newcommand\csname fisica@unidade@G\endcsname{\unit{\cubic\metre\per\kilogram\per\square\second}}
\expandafter\newcommand\csname fisica@simbolo@G\endcsname{$G$}
\expandafter\newcommand\csname fisica@nome@G\endcsname{constante gravitacional}
\expandafter\newcommand\csname fisica@cte@NA\endcsname{$N_A = \qty{6.02214076e23}{\per\mole}$}
\expandafter\newcommand\csname fisica@valor@NA\endcsname{\qty{6.02214076e23}{\per\mole}}
\expandafter\newcommand\csname fisica@unidade@NA\endcsname{\unit{\per\mole}}
\expandafter\newcommand\csname fisica@simbolo@NA\endcsname{$N_A$}
\expandafter\newcommand\csname fisica@nome@NA\endcsname{constante de Avogadro}
\expandafter\newcommand\csname fisica@cte@c\endcsname{$c = \qty{299792458}{\metre\per\second}$}
\expandafter\newcommand\csname fisica@valor@c\endcsname{\qty{299792458}{\metre\per\second}}
\expandafter\newcommand\csname fisica@unidade@c\endcsname{\unit{\metre\per\second}}
\expandafter\newcommand\csname fisica@simbolo@c\endcsname{$c$}
\expandafter\newcommand\csname fisica@nome@c\endcsname{velocidade da luz no vácuo}
\expandafter\newcommand\csname fisica@cte@e\endcsname{$e = \qty{1.602176634e-19}{\coulomb}$}
\expandafter\newcommand\csname fisica@valor@e\endcsname{\qty{1.602176634e-19}{\coulomb}}
\expandafter\newcommand\csname fisica@unidade@e\endcsname{\unit{\coulomb}}
\expandafter\newcommand\csname fisica@simbolo@e\endcsname{$e$}
\expandafter\newcommand\csname fisica@nome@e\endcsname{carga elementar}
\expandafter\newcommand\csname fisica@cte@eps0\endcsname{$\varepsilon_0 = \qty{8.8541878188(14)e-12}{\farad\per\metre}$}
\expandafter\newcommand\csname fisica@valor@eps0\endcsname{\qty{8.8541878188(14)e-12}{\farad\per\metre}}
\expandafter\newcommand\csname fisica@unidade@eps0\endcsname{\unit{\farad\per\metre}}
\expandafter\newcommand\csname fisica@simbolo@eps0\endcsname{$\varepsilon_0$}
\expandafter\newcommand\csname fisica@nome@eps0\endcsname{permissividade elétrica do vácuo}
\expandafter\newcommand\csname fisica@cte@h\endcsname{$h = \qty{6.62607015e-34}{\joule\second}$}
\expandafter\newcommand\csname fisica@valor@h\endcsname{\qty{6.62607015e-34}{\joule\second}}
\expandafter\newcommand\csname fisica@unidade@h\endcsname{\unit{\joule\second}}
\expandafter\newcommand\csname fisica@simbolo@h\endcsname{$h$}
\expandafter\newcommand\csname fisica@nome@h\endcsname{constante de Planck}
\expandafter\newcommand\csname fisica@cte@hbar\endcsname{$\hbar = \qty{1.054571817e-34}{\joule\second}$}
\expandafter\newcommand\csname fisica@valor@hbar\endcsname{\qty{1.054571817e-34}{\joule\second}}
\expandafter\newcommand\csname fisica@unidade@hbar\endcsname{\unit{\joule\second}}
\expandafter\newcommand\csname fisica@simbolo@hbar\endcsname{$\hbar$}
\expandafter\newcommand\csname fisica@nome@hbar\endcsname{constante de Planck reduzida}
\expandafter\newcommand\csname fisica@cte@kB\endcsname{$k_B = \qty{1.380649e-23}{\joule\per\kelvin}$}
\expandafter\newcommand\csname fisica@valor@kB\endcsname{\qty{1.380649e-23}{\joule\per\kelvin}}
\expandafter\newcommand\csname fisica@unidade@kB\endcsname{\unit{\joule\per\kelvin}}
\expandafter\newcommand\csname fisica@simbolo@kB\endcsname{$k_B$}
\expandafter\newcommand\csname fisica@nome@kB\endcsname{constante de Boltzmann}
\expandafter\newcommand\csname fisica@cte@me\endcsname{$m_e = \qty{9.1093837139(28)e-31}{\kilogram}$}
\expandafter\newcommand\csname fisica@valor@me\endcsname{\qty{9.1093837139(28)e-31}{\kilogram}}
\expandafter\newcommand\csname fisica@unidade@me\endcsname{\unit{\kilogram}}
\expandafter\newcommand\csname fisica@simbolo@me\endcsname{$m_e$}
\expandafter\newcommand\csname fisica@nome@me\endcsname{massa do elétron}
\expandafter\newcommand\csname fisica@cte@mp\endcsname{$m_p = \qty{1.67262192595(52)e-27}{\kilogram}$}
\expandafter\newcommand\csname fisica@valor@mp\endcsname{\qty{1.67262192595(52)e-27}{\kilogram}}
\expandafter\newcommand\csname fisica@unidade@mp\endcsname{\unit{\kilogram}}
\expandafter\newcommand\csname fisica@simbolo@mp\endcsname{$m_p$}
\expandafter\newcommand\csname fisica@nome@mp\endcsname{massa do próton}
\expandafter\newcommand\csname fisica@cte@muB\endcsname{$\mu_B = \qty{9.2740100657(29)e-24}{\joule\per\tesla}$}
\expandafter\newcommand\csname fisica@valor@muB\endcsname{\qty{9.2740100657(29)e-24}{\joule\per\tesla}}
\expandafter\newcommand\csname fisica@unidade@muB\endcsname{\unit{\joule\per\tesla}}
\expandafter\newcommand\csname fisica@simbolo@muB\endcsname{$\mu_B$}
\expandafter\newcommand\csname fisica@nome@muB\endcsname{magnéton de Bohr}

\newcommand{\cte}[1]{\csname fisica@cte@#1\endcsname}
\newcommand{\ctevalor}[1]{\csname fisica@valor@#1\endcsname}
\newcommand{\cteunidade}[1]{\csname fisica@unidade@#1\endcsname}
\newcommand{\ctesimbolo}[1]{\csname fisica@simbolo@#1\endcsname}
\newcommand{\tabelaconstantes}[1]{%
  \def\fisica@linhas{}%
  \@for\fisica@chave:=#1\do{%
    \edef\fisica@linhas{%
      \unexpanded\expandafter{\fisica@linhas}%
      \expandafter\noexpand\csname fisica@simbolo@\fisica@chave\endcsname &
      \expandafter\noexpand\csname fisica@valor@\fisica@chave\endcsname &
      \expandafter\noexpand\csname fisica@nome@\fisica@chave\endcsname
      \noexpand\\%
    }%
  }%
  \fisica@linhas
}

\newcommand{\roteiro}{\item \textbf{Motivação} --- Por que o oscilador harmônico aparece em toda a física. \item \textbf{Formulação quântica} --- A equação de Schrödinger e a quantização da energia. \item \textbf{Espectro e funções de onda} --- Níveis, paridade, nós e densidade de probabilidade. \item \textbf{Aplicações e ferramentas} --- Do infravermelho aos qubits; constantes e dados via Lua. }
\newcommand{\comparacao}{Energia & contínua, $E \ge 0$ & discreta, $E_n=(n+\tfrac12)\hbar\omega$ \\Estado de menor energia & repouso no fundo do poço, $E=0$ & energia de ponto zero, $E_0=\tfrac12\hbar\omega$ \\Posição & trajetória $x(t)$ bem definida & distribuição de probabilidade $|\psi_n|^2$ \\Amplitude & qualquer valor permitido por $E$ & probabilidade não nula fora da região clássica \\Radiação emitida & qualquer frequência & saltos discretos $\Delta E=\hbar\omega$ \\}
\newcommand{\niveisTabela}{0 & \qty{125}{\milli\electronvolt} & $+1$ & 0 \\1 & \qty{375}{\milli\electronvolt} & $-1$ & 1 \\2 & \qty{625}{\milli\electronvolt} & $+1$ & 2 \\3 & \qty{875}{\milli\electronvolt} & $-1$ & 3 \\4 & \qty{1125}{\milli\electronvolt} & $+1$ & 4 \\5 & \qty{1375}{\milli\electronvolt} & $-1$ & 5 \\}
\newcommand{\exercicios}{\item Partindo de $[\hat x,\hat p]=i\hbar$, mostre que $[\hat H,\hat a^\dagger]=\hbar\omega\,\hat a^\dagger$.\hfill{\scriptsize\color{fisicamuted}nível 1} \item Calcule $\langle x\rangle$ e $\langle x^2\rangle$ no estado fundamental e verifique o princípio da incerteza.\hfill{\scriptsize\color{fisicamuted}nível 2} \item Determine a probabilidade de encontrar a partícula fora da região classicamente permitida para $n=0$.\hfill{\scriptsize\color{fisicamuted}nível 3} \item Estime $E_0$ por análise dimensional e compare com $\tfrac12\hbar\omega$.\hfill{\scriptsize\color{fisicamuted}nível 1} \item Esboce $\psi_3(\xi)$ e identifique seus nós e sua paridade.\hfill{\scriptsize\color{fisicamuted}nível 1} }

\newcommand{\plotOscilador}{\addplot[draw=fisicamuted, dashed, line width=0.7pt, smooth] coordinates {(-4.0000,8.0000) (-3.9200,7.6832) (-3.8400,7.3728) (-3.7600,7.0688) (-3.6800,6.7712) (-3.6000,6.4800) (-3.5200,6.1952) (-3.4400,5.9168) (-3.3600,5.6448) (-3.2800,5.3792) (-3.2000,5.1200) (-3.1200,4.8672) (-3.0400,4.6208) (-2.9600,4.3808) (-2.8800,4.1472) (-2.8000,3.9200) (-2.7200,3.6992) (-2.6400,3.4848) (-2.5600,3.2768) (-2.4800,3.0752) (-2.4000,2.8800) (-2.3200,2.6912) (-2.2400,2.5088) (-2.1600,2.3328) (-2.0800,2.1632) (-2.0000,2.0000) (-1.9200,1.8432) (-1.8400,1.6928) (-1.7600,1.5488) (-1.6800,1.4112) (-1.6000,1.2800) (-1.5200,1.1552) (-1.4400,1.0368) (-1.3600,0.9248) (-1.2800,0.8192) (-1.2000,0.7200) (-1.1200,0.6272) (-1.0400,0.5408) (-0.9600,0.4608) (-0.8800,0.3872) (-0.8000,0.3200) (-0.7200,0.2592) (-0.6400,0.2048) (-0.5600,0.1568) (-0.4800,0.1152) (-0.4000,0.0800) (-0.3200,0.0512) (-0.2400,0.0288) (-0.1600,0.0128) (-0.0800,0.0032) (0.0000,0.0000) (0.0800,0.0032) (0.1600,0.0128) (0.2400,0.0288) (0.3200,0.0512) (0.4000,0.0800) (0.4800,0.1152) (0.5600,0.1568) (0.6400,0.2048) (0.7200,0.2592) (0.8000,0.3200) (0.8800,0.3872) (0.9600,0.4608) (1.0400,0.5408) (1.1200,0.6272) (1.2000,0.7200) (1.2800,0.8192) (1.3600,0.9248) (1.4400,1.0368) (1.5200,1.1552) (1.6000,1.2800) (1.6800,1.4112) (1.7600,1.5488) (1.8400,1.6928) (1.9200,1.8432) (2.0000,2.0000) (2.0800,2.1632) (2.1600,2.3328) (2.2400,2.5088) (2.3200,2.6912) (2.4000,2.8800) (2.4800,3.0752) (2.5600,3.2768) (2.6400,3.4848) (2.7200,3.6992) (2.8000,3.9200) (2.8800,4.1472) (2.9600,4.3808) (3.0400,4.6208) (3.1200,4.8672) (3.2000,5.1200) (3.2800,5.3792) (3.3600,5.6448) (3.4400,5.9168) (3.5200,6.1952) (3.6000,6.4800) (3.6800,6.7712) (3.7600,7.0688) (3.8400,7.3728) (3.9200,7.6832) (4.0000,8.0000)}; \addplot[draw=fisicacyan, line width=1.5pt, smooth] coordinates {(-4.0000,0.5001) (-3.9200,0.5001) (-3.8400,0.5001) (-3.7600,0.5002) (-3.6800,0.5003) (-3.6000,0.5003) (-3.5200,0.5005) (-3.4400,0.5006) (-3.3600,0.5008) (-3.2800,0.5010) (-3.2000,0.5013) (-3.1200,0.5017) (-3.0400,0.5022) (-2.9600,0.5028) (-2.8800,0.5036) (-2.8000,0.5045) (-2.7200,0.5056) (-2.6400,0.5069) (-2.5600,0.5085) (-2.4800,0.5104) (-2.4000,0.5126) (-2.3200,0.5153) (-2.2400,0.5183) (-2.1600,0.5219) (-2.0800,0.5259) (-2.0000,0.5305) (-1.9200,0.5357) (-1.8400,0.5415) (-1.7600,0.5479) (-1.6800,0.5549) (-1.6000,0.5627) (-1.5200,0.5710) (-1.4400,0.5799) (-1.3600,0.5894) (-1.2800,0.5993) (-1.2000,0.6097) (-1.1200,0.6203) (-1.0400,0.6312) (-0.9600,0.6421) (-0.8800,0.6530) (-0.8000,0.6636) (-0.7200,0.6739) (-0.6400,0.6836) (-0.5600,0.6926) (-0.4800,0.7008) (-0.4000,0.7080) (-0.3200,0.7141) (-0.2400,0.7189) (-0.1600,0.7225) (-0.0800,0.7246) (0.0000,0.7253) (0.0800,0.7246) (0.1600,0.7225) (0.2400,0.7189) (0.3200,0.7141) (0.4000,0.7080) (0.4800,0.7008) (0.5600,0.6926) (0.6400,0.6836) (0.7200,0.6739) (0.8000,0.6636) (0.8800,0.6530) (0.9600,0.6421) (1.0400,0.6312) (1.1200,0.6203) (1.2000,0.6097) (1.2800,0.5993) (1.3600,0.5894) (1.4400,0.5799) (1.5200,0.5710) (1.6000,0.5627) (1.6800,0.5549) (1.7600,0.5479) (1.8400,0.5415) (1.9200,0.5357) (2.0000,0.5305) (2.0800,0.5259) (2.1600,0.5219) (2.2400,0.5183) (2.3200,0.5153) (2.4000,0.5126) (2.4800,0.5104) (2.5600,0.5085) (2.6400,0.5069) (2.7200,0.5056) (2.8000,0.5045) (2.8800,0.5036) (2.9600,0.5028) (3.0400,0.5022) (3.1200,0.5017) (3.2000,0.5013) (3.2800,0.5010) (3.3600,0.5008) (3.4400,0.5006) (3.5200,0.5005) (3.6000,0.5003) (3.6800,0.5003) (3.7600,0.5002) (3.8400,0.5001) (3.9200,0.5001) (4.0000,0.5001)}; \addlegendentry{$n=0$} \addplot[draw=fisicaamber, line width=1.5pt, smooth] coordinates {(-4.0000,1.4996) (-3.9200,1.4994) (-3.8400,1.4992) (-3.7600,1.4990) (-3.6800,1.4987) (-3.6000,1.4982) (-3.5200,1.4977) (-3.4400,1.4970) (-3.3600,1.4962) (-3.2800,1.4952) (-3.2000,1.4939) (-3.1200,1.4923) (-3.0400,1.4905) (-2.9600,1.4882) (-2.8800,1.4855) (-2.8000,1.4823) (-2.7200,1.4786) (-2.6400,1.4742) (-2.5600,1.4692) (-2.4800,1.4635) (-2.4000,1.4571) (-2.3200,1.4499) (-2.2400,1.4419) (-2.1600,1.4332) (-2.0800,1.4238) (-2.0000,1.4137) (-1.9200,1.4031) (-1.8400,1.3921) (-1.7600,1.3808) (-1.6800,1.3694) (-1.6000,1.3582) (-1.5200,1.3474) (-1.4400,1.3373) (-1.3600,1.3281) (-1.2800,1.3202) (-1.2000,1.3139) (-1.1200,1.3094) (-1.0400,1.3070) (-0.9600,1.3070) (-0.8800,1.3096) (-0.8000,1.3149) (-0.7200,1.3229) (-0.6400,1.3338) (-0.5600,1.3474) (-0.4800,1.3637) (-0.4000,1.3823) (-0.3200,1.4031) (-0.2400,1.4257) (-0.1600,1.4497) (-0.0800,1.4746) (0.0000,1.5000) (0.0800,1.5254) (0.1600,1.5503) (0.2400,1.5743) (0.3200,1.5969) (0.4000,1.6177) (0.4800,1.6363) (0.5600,1.6526) (0.6400,1.6662) (0.7200,1.6771) (0.8000,1.6851) (0.8800,1.6904) (0.9600,1.6930) (1.0400,1.6930) (1.1200,1.6906) (1.2000,1.6861) (1.2800,1.6798) (1.3600,1.6719) (1.4400,1.6627) (1.5200,1.6526) (1.6000,1.6418) (1.6800,1.6306) (1.7600,1.6192) (1.8400,1.6079) (1.9200,1.5969) (2.0000,1.5863) (2.0800,1.5762) (2.1600,1.5668) (2.2400,1.5581) (2.3200,1.5501) (2.4000,1.5429) (2.4800,1.5365) (2.5600,1.5308) (2.6400,1.5258) (2.7200,1.5214) (2.8000,1.5177) (2.8800,1.5145) (2.9600,1.5118) (3.0400,1.5095) (3.1200,1.5077) (3.2000,1.5061) (3.2800,1.5048) (3.3600,1.5038) (3.4400,1.5030) (3.5200,1.5023) (3.6000,1.5018) (3.6800,1.5013) (3.7600,1.5010) (3.8400,1.5008) (3.9200,1.5006) (4.0000,1.5004)}; \addlegendentry{$n=1$} \addplot[draw=fisicaviolet, line width=1.5pt, smooth] coordinates {(-4.0000,2.5017) (-3.9200,2.5022) (-3.8400,2.5029) (-3.7600,2.5037) (-3.6800,2.5048) (-3.6000,2.5061) (-3.5200,2.5077) (-3.4400,2.5097) (-3.3600,2.5122) (-3.2800,2.5151) (-3.2000,2.5185) (-3.1200,2.5226) (-3.0400,2.5274) (-2.9600,2.5330) (-2.8800,2.5393) (-2.8000,2.5464) (-2.7200,2.5544) (-2.6400,2.5632) (-2.5600,2.5728) (-2.4800,2.5832) (-2.4000,2.5941) (-2.3200,2.6055) (-2.2400,2.6171) (-2.1600,2.6288) (-2.0800,2.6402) (-2.0000,2.6509) (-1.9200,2.6608) (-1.8400,2.6692) (-1.7600,2.6759) (-1.6800,2.6805) (-1.6000,2.6825) (-1.5200,2.6817) (-1.4400,2.6778) (-1.3600,2.6706) (-1.2800,2.6599) (-1.2000,2.6458) (-1.1200,2.6284) (-1.0400,2.6079) (-0.9600,2.5847) (-0.8800,2.5594) (-0.8000,2.5324) (-0.7200,2.5045) (-0.6400,2.4765) (-0.5600,2.4492) (-0.4800,2.4234) (-0.4000,2.4000) (-0.3200,2.3796) (-0.2400,2.3630) (-0.1600,2.3507) (-0.0800,2.3432) (0.0000,2.3407) (0.0800,2.3432) (0.1600,2.3507) (0.2400,2.3630) (0.3200,2.3796) (0.4000,2.4000) (0.4800,2.4234) (0.5600,2.4492) (0.6400,2.4765) (0.7200,2.5045) (0.8000,2.5324) (0.8800,2.5594) (0.9600,2.5847) (1.0400,2.6079) (1.1200,2.6284) (1.2000,2.6458) (1.2800,2.6599) (1.3600,2.6706) (1.4400,2.6778) (1.5200,2.6817) (1.6000,2.6825) (1.6800,2.6805) (1.7600,2.6759) (1.8400,2.6692) (1.9200,2.6608) (2.0000,2.6509) (2.0800,2.6402) (2.1600,2.6288) (2.2400,2.6171) (2.3200,2.6055) (2.4000,2.5941) (2.4800,2.5832) (2.5600,2.5728) (2.6400,2.5632) (2.7200,2.5544) (2.8000,2.5464) (2.8800,2.5393) (2.9600,2.5330) (3.0400,2.5274) (3.1200,2.5226) (3.2000,2.5185) (3.2800,2.5151) (3.3600,2.5122) (3.4400,2.5097) (3.5200,2.5077) (3.6000,2.5061) (3.6800,2.5048) (3.7600,2.5037) (3.8400,2.5029) (3.9200,2.5022) (4.0000,2.5017)}; \addlegendentry{$n=2$} }

\newcommand{\energia}[2]{%
  \qty[drop-zero-decimal]{\fpeval{round((#1+0.5)*(#2),3)}}{\milli\electronvolt}%
}

\makeatother
